\documentclass[11pt]{article}
\usepackage[utf8]{inputenc}
\usepackage{amsmath,amssymb,amsfonts}
\usepackage{booktabs}
\usepackage{graphicx}
\usepackage{hyperref}
\usepackage{geometry}
\geometry{margin=1in}

\title{StateShift: Tracking State-Dependent Reasoning Recovery Across Post-Training}
\author{Anonymous Authors}
\date{}

\begin{document}
\maketitle

\begin{abstract}
Standard evaluations of reinforcement learning (RL) post-training in large language models rely primarily on aggregate benchmark accuracy, masking how local reasoning capabilities evolve. We introduce \textbf{StateShift}, a controlled framework that evaluates target-transition recovery conditional on intermediate reasoning state. In a controlled study ($N=454, 29,056$ rollouts), RL post-training yields an 11.76-percentage-point state-by-checkpoint interaction in target-transition success ($\Gamma_{256} = +0.1176, p < 0.0001, 95\%$ problem-blocked bootstrap CI $[+0.0955, +0.1400]$), demonstrating that post-training gains are state-selective rather than uniform. Tracking this contrast across nine empirically sampled checkpoints ($t \in \{0, 32, 64, 96, 128, 160, 192, 224, 256\}$) reveals an overall non-decreasing sampled trajectory ($0.0000 \to +0.0333 \to +0.0337 \to +0.0774 \to +0.0748 \to +0.0598 \to +0.0976 \to +0.0950 \to +0.1176$), with state-selective interaction statistically detectable by the earliest available post-training checkpoint ($t=32, \Gamma_{32} = +0.0333$, Bonferroni 95\% CI $[+0.0011, +0.0655]$). Finally, evaluating unperturbed rollouts ($N=200, 3,200$ rollouts) reveals a $18.19\%$ Natural Error Incidence ($\text{NEI}$) and a $30.93\%$ Conditional Natural Post-Error Recovery Rate ($\text{NRR}$, 95\% CI $[27.19\%, 34.82\%]$) after verifier-confirmed natural reasoning errors. Our results demonstrate that state-conditioned evaluation reveals fine-grained behavioral dynamics obscured by aggregate correctness metrics alone.
\end{abstract}

\section{Introduction}
Reinforcement learning from verifier feedback (RLVR) has significantly improved mathematical reasoning in large language models. However, standard benchmark metrics measure only final accuracy, providing limited insight into how post-training changes reasoning behavior conditional on local state. Does RL training uniformly improve reasoning capability, or does it specifically enhance the model's ability to recover from invalid intermediate states?

To answer this question, we present \textbf{StateShift}, a rigorous experimental framework evaluating state-conditioned reasoning dynamics across post-training.

\section{The StateShift Framework}
StateShift evaluates target-transition success under two matched conditions:
\begin{enumerate}
    \item \textbf{Recovery Condition ($R$)}: The model is initialized at a locally invalid intermediate reasoning state.
    \item \textbf{Control Condition ($C$)}: The model is initialized at a matched control state.
\end{enumerate}

The primary estimand is the difference-in-differences state-by-checkpoint interaction:
\begin{equation}
\Gamma_t = (\mu_{R,t} - \mu_{R,0}) - (\mu_{C,t} - \mu_{C,0})
\end{equation}

\section{Controlled Endpoint Experiment (Study A)}
In our primary confirmatory experiment ($N=454, K=16, 29,056$ rollouts), we observe:
\begin{align}
\mu_{R,0} = 0.3834, \quad \mu_{C,0} = 0.3892 \\
\mu_{R,256} = 0.7039, \quad \mu_{C,256} = 0.5921
\end{align}
yielding a primary interaction $\Gamma_{256} = \mathbf{+0.1176}$ ($95\%$ problem-blocked bootstrap CI $[+0.0955, +0.1400], p < 0.0001$). Under strict decontamination filtering ($N_{\text{Strict}}=388$), $\Gamma_{256,\text{Strict}} = \mathbf{+0.1160}$ ($95\%$ CI $[+0.0913, +0.1408]$).

\section{Complete Nine-Checkpoint Empirical Trajectory}
Evaluating intermediate checkpoints across training yields the empirical nine-point vector:
\begin{equation}
\mathbf{\Gamma} = [0.0000, +0.0333, +0.0337, +0.0774, +0.0748, +0.0598, +0.0976, +0.0950, +0.1176]
\end{equation}
An order-restricted test confirms a non-decreasing overall sampled trajectory. The positive interaction is statistically detectable by the earliest available checkpoint $t=32$ ($\Gamma_{32} = +0.0333$, Bonferroni 95\% CI $[+0.0011, +0.0655]$).

\section{Natural Post-Error Recovery (Study B)}
In $3,200$ unperturbed rollouts ($N=200, K=16$), the model exhibits a Natural Error Incidence $\text{NEI} = 18.19\%$ ($582$ error rollouts). Conditional on a verifier-confirmed natural reasoning error, the model autonomously recovers to the correct final answer in $180$ of $582$ qualifying episodes ($\text{NRR} = \mathbf{30.93\%}$, $95\%$ problem-blocked bootstrap CI $[27.19\%, 34.82\%]$).

\section{Limitations}
Strict pairwise monotonicity across all adjacent intervals was not established due to sample size constraints ($K=2/3$ at intermediate steps). Sub-32-step emergence remains unobserved because fine-grained checkpoints $t \in \{1..31\}$ do not exist. Natural post-error recovery is conditional on error occurrence and should not be interpreted as an unqualified causal self-correction probability.

\section{Conclusion}
State-conditioned evaluation provides a powerful tool for analyzing LLM post-training dynamics, proving that RL induces a state-selective transition recovery capability that operates in both controlled perturbations and unprompted natural decoding.

\bibliographystyle{plain}
\bibliography{references}

\end{document}
